\chapter{Conclusion and Future Work}
\label{chap:conclusion}

\section{Summary}
This thesis presented \textbf{EdgeVerify}, a Joint-Embedding Predictive
Architecture designed for localized execution on edge and Dew computing
nodes. The architecture combines an EMA-updated target encoder, an
action-conditioned latent predictor, and a VICReg-inspired
variance/covariance objective applied within a local, layer-wise
distillation loss. By predicting in an abstract latent space rather than
reconstructing raw inputs, and by bounding the backpropagation graph
through stop-gradient firebreaks, the design targets the two coupled
obstacles to on-device learning identified in Chapter~\ref{chap:intro}:
strict memory budgets and representational collapse.

\section{Contributions}
The work contributes: (i) the EdgeVerify architecture and its
action-conditioned prediction formulation; (ii) a negative-free, stable
local loss that regularizes both the encoder output (variance) and the
predictor output (covariance); and (iii) a proof-of-concept PyTorch
implementation together with a reproducible protocol for evaluating
convergence, collapse, latency, and memory footprint.

\section{Limitations}
The principal limitation of the present work is that empirical validation
is preliminary. The evaluation in Chapter~\ref{chap:results} was conducted
on non-standard data (procedurally generated structured images and the
low-resolution scikit-learn digits set), using a single random seed and
CPU-only measurement, because the environment lacked network access to
standard benchmarks. Within those limits, the evaluation surfaced a
substantive finding: while the variance term prevents variance collapse,
the covariance term at its current weight ($\beta = 0.01$) fails to prevent
dimensional (rank) collapse, and the unregularized baseline was in fact
higher-rank. The design's collapse-prevention claim is therefore only
partially supported and requires the regularizer-weight ablation described
below. In addition, the
convolutional encoder used in the proof-of-concept is small relative to the
transformer backbones typical of the JEPA literature, so conclusions drawn
here may not transfer directly to larger models. The ``multi-agent'' aspect
suggested by the title is, at present, realized only as the multi-module
(context/target/predictor) structure of a single node; coordination across
multiple physical agents is not yet implemented.

\section{Future Work}
Several directions follow directly:
\begin{itemize}
    \item \textbf{Empirical evaluation.} Implement the dataset loader
    (\texttt{dataset.py}) and benchmark suite (\texttt{evaluate.py}) and
    run the protocol of Chapter~\ref{chap:results} on a standard vision
    benchmark and a defined edge device, populating the results tables.
    \item \textbf{Regularizer-weight ablation.} Systematically increase the
    covariance weight $\beta$ (and vary $\alpha$ and the threshold $\gamma$)
    to test whether stronger decorrelation restores effective rank, and
    evaluate applying the variance and covariance penalties to the same
    embedding rather than to $s_t$ and $s_{\text{pred}}$ separately.
    \item \textbf{Backbone scaling.} Evaluate whether the local objective
    remains stable with larger or transformer-based encoders.
    \item \textbf{True multi-agent extension.} Extend the single-node
    design to coordinated learning across multiple edge nodes, connecting
    the framework to decentralized-learning approaches such as federated
    learning \cite{mcmahan2017communication}.
\end{itemize}
